\section{Implementation Defects and Corrections}
\label{app:defects}
%==============================================================================

This appendix records defects found during development, their effect on reported numbers, and their corrections. Every table in the paper reflects the corrected code.

\subsection*{Loop detector and reproducibility}
\begin{itemize}[leftmargin=*]
  \item \textbf{Window-sampling seed.} The calibration and validation script seeded its window sampling with Python's built-in \texttt{hash()} on corpus names, which is randomized per process, so the sampled windows differed between invocations. It now uses a seed derived from \texttt{hashlib.md5}, verified identical across processes, and saves the sampled window start indices. Re-running reproduces the frozen setting ($R{=}5$, $P_{\max}{=}60$) and closely similar false-positive counts.
  \item \textbf{Period-1 threshold.} The detector was first described as ``$R{=}5$'' without qualification; period 1 uses a decoupled threshold of 8.
  \item \textbf{Result keys and prompts.} \texttt{compute\_\allowbreak{}significance.py} used the pre-rename key \texttt{lt\_lookback\_a3.0}, which silently dropped that comparison from every run after the class was renamed from \texttt{LookBackDecoder} (a name that implied the published Look-back algorithm) to \textsc{SuffixMatchDecoder}. Greedy decoding was initially evaluated on 5 of the 15 test prompts and is now evaluated on all 15.
\end{itemize}

\subsection*{Model-consistency NLL}
The first implementation scored non-overlapping chunks of length \texttt{block\_\allowbreak{}size}; on a toy model about half of all scored positions received less context than the architecture allows (in one case 9 of 16 possible context tokens). A first fix (overlapping windows, keeping each window's back half) was still insufficient. The final method scores every target position independently with its exact context window; on the toy model all 40 scored positions receive exactly $\min(t,\text{block\_size})$ tokens. An intermediate version left-padded windows with a real vocabulary token id, which unmasked causal attention attends to (0.26 maximum logit difference between isolated and padded evaluation of the same context); right-padding is used instead, since causal masking cannot let it leak into real positions, and the batched computation matches an isolated forward pass to 6 decimals.

\subsection*{Dual solver}
\begin{enumerate}[leftmargin=*]
  \item \textbf{Float32 underflow.} The solver computed $q_\lambda$'s logits as $\log(p+10^{-40})-\lambda\risk$ from a materialized softmax. At $T{=}0.10$, 276 of 2000 tested vocabulary entries underflowed to exactly $0.0$ in one constructed check, with true log-probabilities spanning 41 nats collapsed to one floor value ($\log10^{-40}\approx-92.1$). The solver now works in log-space throughout.
  \item \textbf{Infeasibility conflated with the cap.} True infeasibility ($\varepsilon<\min_v\risk(v)$) was learned only by bisecting to a fixed cap ($\lambda_{\max}{=}50$, later 1000) and failing. It is now decided in closed form and the bracket expands geometrically (at most 60 doublings; an earlier description of the search as having ``no ceiling'' was incorrect).
  \item \textbf{Discarded diagnostics.} An earlier run hid a solver that was silently infeasible at about 26--28\% of steps.
  \item \textbf{Float32/float64 comparison at the boundary.} At $\varepsilon=\min_v\risk(v)$ the comparison could round the wrong way and misreport a feasible point as infeasible; a tolerance and a \texttt{near\_\allowbreak{}boundary} flag were added.
  \item \textbf{NaN diagnostics.} A NaN entered the KL and entropy diagnostics whenever top-$p$ filtering ran after the projection with $q(v){=}0$ exactly.
\end{enumerate}
\textbf{Effect.} Dual per-step KL on the primary model was 2.256--2.328 bits in the original run, 10.3--11.6 bits after the cap fix alone, and 11.2--12.7 bits after the numerical fix, about five times the original figure. Per-step feasibility was 72.3--74.5\% originally, 98.75--98.87\% after the cap fix, and is exactly 100.0\% after the numerical fix. Loop counts do not depend on this arithmetic and were unaffected.

\subsection*{Theory}
The projection theorem originally omitted the $\varepsilon=\min_v\risk(v)$ boundary, where no finite dual variable exists. The statement that the method reduces to hard no-repeat blocking as $\lambda\to\infty$ was wrong twice: it fails when $n_{\min}<n_{\max}$, and even for a single order it ignores that hard blocking fails when no safe continuation exists. A statement that no causally committed token-by-token rule can realize the event-level optimum was wrong: the affine-hazard kernel (Corollary~\ref{cor:affine}) is such a rule when the exact event-to-go hazard is available. The no-repeat argument previously stated a lag condition $P\le n$ that is not needed. And an earlier explanation of the gap in the small system attributed it to using the hazard exponentially instead of affinely; the ablation of Section~\ref{ssec:gap} shows that the functional form accounts for at most about 10\% of the optimum's KL and that re-solving a fixed budget at every step accounts for nearly all of the gap.

\subsection*{Small-system experiments}
\begin{itemize}[leftmargin=*]
  \item \textbf{Unreproducible headline.} An earlier version reported ``17.9\% vs.\ 53.2\% loop rate at matched KL$\approx0.24$'', which could not be reproduced from its saved artifact: the saved check targeted a different operating point than the one cited. The comparison is now computed and saved for every grid point.
  \item \textbf{Mismatched $\varepsilon$.} Local and hazard-informed control were called with a one-step $\varepsilon$ read directly off the grid used for the optimum's sequence-level target, so the mechanisms were compared at different operating points. All are now matched at the same sequence-level loop probability.
  \item \textbf{Ignored infeasibility flag.} The one-step solver's structural-infeasibility flag was computed but never checked; it is now tracked and reported.
  \item \textbf{Non-absorbing KL.} The Monte Carlo simulator kept charging KL for the full horizon after a trajectory had looped; KL is now absorbing at the loop, matching the exact derivation.
  \item \textbf{Approximate matching.} The equal-distortion table used a Monte Carlo simulator (2000 trajectories per point) and the closest point on a 60-point grid of per-step budgets, so its matched KL differed from the target in the third decimal and its loop probabilities carried sampling noise. The simulator is removed and the budgets are solved exactly; the exact local loop probabilities (44.5, 45.2, 42.1, 25.4, 7.4\%) replace the earlier simulated ones (43.9, 44.2, 41.1, 24.9, 9.6\%), and the exact KL ceiling is $0.614985$ bits per step instead of an estimated $\approx0.613$.
  \item \textbf{Base of the exponent.} The derivation mixed $2^{-\lambda}$ and $e^{-\lambda}$; it uses base 2 throughout, and the closed-form identity holds to $<10^{-14}$.

\item \textbf{Infeasibility confounded with allocation.} The first small-system comparison attributed the excess KL of the per-step controllers to local control in general and did not separate two effects: the hazard-informed controller's budget was below its achievable floor at 33--50\% of the probability mass (it then collapses onto the safest token), and a fixed budget allocates distortion across histories badly even where it is feasible. Section~\ref{sec:small} separates them with an always-feasible budget, an abstaining variant, a system with no infeasible node at all, and the second-order waste law.
\end{itemize}

\subsection*{Baseline reimplementations}
\textbf{FSD-style.} Three versions were implemented. The first combined the model and the anti-LM in logit space with a longest-available-order heuristic over the full vocabulary. The second used the paper's probability-space scoring over the top-$k$ candidates, with fixed normalized interpolation weights $\beta^{N-k}$ (about $0.30/0.33/0.37$ for orders $1/2/3$ at $N{=}3$, $\beta{=}0.9$) and $\text{count}/\text{total}$ at every order. The official anti-LM instead allocates residual mass recursively and smooths higher orders by $\text{count}/(\text{total}{+}1)$ (Appendix~\ref{app:fsd}). The third version, checked against the official function, is the one reported. Between the second and the third version the primary-model loop count moved from 2/150 to 3/150, non-fiction from 16.7\% to 16.0\%, GPT-2 mcNLL from 2.122 to 2.178 (0.0\% loops in both), OLMo KL from 3.188 to 3.194 bits, Mistral loops from 0/50 to 1/50 and KL from 4.252 to 4.459 bits, common-KL KL from 2.184 to 2.148 bits, and the spelling-error rate from 2.59\% to 2.29\%; the first-to-second version change moved the primary-model loop rate from 0.0\% to 1.3\%.

\textbf{LZ Penalty.} An earlier version scored a single identified match instead of every candidate, used a different formula for non-matching tokens, and had the opposite sign convention. The ported version, checked against the reference at the pinned commit, moved the primary-model KL to 0.186 bits, non-fiction loops from 0.0\% to 2.0\%, GPT-2 loops from 25.0\% to 8.0\%, Mistral loops from 6.0\% to 2.0\%, and the spelling-error rate from 1.58\% to 0.75\%. Its history handling still differs from the reference (Appendix~\ref{app:lz}).

\subsection*{Real-model KL analysis}
An earlier version of Section~\ref{ssec:eventlevel} pooled all prompts into a single reference loop rate and reported that every decoder spends ``at least 50 times'' the event-level term. Pooling prompts lowers the bound relative to the prompt-conditional one (joint convexity of $d_{\mathrm{Ber}}$) and treats estimated quantities as exact, so it overstated how much distortion could be avoided. The analysis now conditions on the prompt, keeps event counts and path KL together, resamples whole prompts, and reports plug-in estimates with intervals.

\subsection*{Hazard estimator}
The first rollout estimator (i) detected the loop on the continuation alone, discarding the earlier context, so that a cycle begun in the context and completed in the rollout was missed, which is the case a hazard estimate must capture; and (ii) left-padded variable-length contexts with a real token id and no attention mask, which changes the next-token distribution of the model. Its earlier null result (mean hazard 0.027 for risky and 0.028 for novel candidates, $n{=}156$ and $1644$, and no separation in the multi-step variant) therefore came from an estimator that did not measure the quantity it was named for and is not used. Both defects are fixed (Appendix~\ref{sec:hazard-transfer}).

\subsection*{Pilot human evaluation and release}
The pilot's FSD-style passages came from the intermediate FSD-style version, which an earlier draft described inconsistently in two places; Appendix~\ref{sec:human-eval} identifies the version by matching the passages against saved runs. A public snapshot of the repository taken before the final changes contained the earlier FSD-style code and lacked the tests and the new result files; the release now contains the code, tests, and results for every table (Appendix~\ref{sec:reproducibility}).

%==============================================================================
